% Appendix: Toy Derivation of \kappa in a Diffusive Substrate Model
% Track P (Core Physics) — Toy Model Demonstration

\section{Appendix: Toy Derivation of $\kappa$ in a Diffusive Substrate Model}
\label{app:kappa-toy}

In this appendix we construct an explicit, checkable toy model in which the
conversion constant $\kappa$ in $\varphi = \kappa \, \tau_{\mathrm{lat}}$
is \emph{not} a free parameter, but is instead determined by substrate-level
quantities.

The goal is not to model the full Leech-lattice substrate $\substrate$, but to
provide a concrete existence proof: there exist reasonable substrate dynamics
for which the interface potential $\varphi$ obeys a Poisson equation and
$\kappa$ is fixed by microscopic parameters.

\subsection{Substrate Geometry and Time Discretization}
We consider a 3D cubic lattice $\mathbb{Z}^3$ with lattice spacing $a$, so that
physical positions are $x = a i$ with $i \in \mathbb{Z}^3$. Time is discrete
in steps of duration $\Delta t$, with integer index $n$ and physical time
$t_n = n\,\Delta t$.

At each site $i$ and time $n$ we define:
\begin{itemize}
  \item $\rho_i^{(n)}$: local processing load / mass density (bits per unit volume $a^3$),
  \item $\tau_i^{(n)}$: mapping latency (cycles per bit) at that site.
\end{itemize}

We also introduce substrate parameters:
\begin{itemize}
  \item $C$: per-site processing capacity in bits per time step $\Delta t$,
  \item $f_U$: Universal Clock Frequency (cycles per second),
  \item $\beta$: coupling constant relating density to bit-arrival rate.
\end{itemize}

\subsection{Operational Definition of Latency}
We model the local mapping latency using a simple queueing-theoretic picture.
Let $\lambda_i^{(n)}$ be the expected number of arriving bits at site $i$
during one time step $\Delta t$. We assume
\begin{equation}
\lambda_i^{(n)} = \beta \, \rho_i^{(n)} a^3,
\end{equation}
so that higher local density induces more arrivals.

We take each site to act as an M/M/1-like queue with service capacity $C$
(bits per time step). In the stable regime $0 < \lambda_i^{(n)} < C$, the mean
waiting time (in units of $\Delta t$) for a bit in a single-server queue is
proportional to $(C - \lambda_i^{(n)})^{-1}$. For concreteness we take
\begin{equation}
T_{\mathrm{wait},i}^{(n)} = \frac{\Delta t}{C - \lambda_i^{(n)}},
\end{equation}
so the latency per bit in cycles/bit is
\begin{equation}
\tau_i^{(n)} = f_U \, T_{\mathrm{wait},i}^{(n)} = \frac{f_U \, \Delta t}{C - \beta \, \rho_i^{(n)} a^3}.
\end{equation}

In the weak-load regime $\beta \, \rho_i^{(n)} a^3 \ll C$, we can expand in a
Taylor series:
\begin{equation}
\tau_i^{(n)} \approx \frac{f_U \Delta t}{C}\left(1 + \frac{\beta a^3}{C} \, \rho_i^{(n)} + O((\rho_i^{(n)})^2)\right).
\end{equation}
We separate a baseline latency $\tau_0$ and a density-dependent perturbation
$\delta \tau_i^{(n)}$:
\begin{align}
\tau_0 &= \frac{f_U \Delta t}{C}, \\
\delta \tau_i^{(n)} &\approx \tau_0 \, \frac{\beta a^3}{C} \, \rho_i^{(n)}.
\end{align}
Thus, to first order, $\tau$ is affine in $\rho$.

\subsection{Spatial Coupling: Discrete Diffusion of Latency}
We now introduce spatial coupling between sites to model the UOS's tendency to
redistribute computational burden and smooth latency differences. Let $N(i)$
denote the six nearest neighbors of $i$ in the cubic lattice. We define the
update rule
\begin{equation}
\tau_i^{(n+1)} - \tau_i^{(n)} = \gamma \sum_{j \in N(i)} \left(\tau_j^{(n)} - \tau_i^{(n)}\right)
  + \alpha \left(\rho_i^{(n)} - \rho_0\right),
\label{eq:discrete-update}
\end{equation}
where
\begin{itemize}
  \item $\gamma$ controls the strength of diffusion-like smoothing,
  \item $\alpha$ couples local density perturbations to latency,
  \item $\rho_0$ is a reference density (vacuum level).
\end{itemize}

The first term drives $\tau$ towards the local average of its neighbors; the
second term encodes the idea that regions with higher density induce larger
latency.

\subsection{Continuum Limit and Effective PDE}
We now pass to a continuum approximation. Let $x_i = a i$ and $t_n = n\Delta t$,
and define fields $\tau_{\mathrm{lat}}(x,t)$ and $\rho(x,t)$ such that
\begin{equation}
\tau_i^{(n)} \approx \tau_{\mathrm{lat}}(x_i, t_n), \quad
\rho_i^{(n)} \approx \rho(x_i, t_n).
\end{equation}

The left-hand side of~\eqref{eq:discrete-update} becomes a finite difference
approximation to a time derivative:
\begin{equation}
\frac{\tau_i^{(n+1)} - \tau_i^{(n)}}{\Delta t} \to \partial_t \tau_{\mathrm{lat}}(x,t)
\end{equation}
as $\Delta t \to 0$.

For the spatial term, using a standard 3D 6-point stencil, we have
\begin{equation}
\sum_{j \in N(i)} (\tau_j^{(n)} - \tau_i^{(n)})
  \approx a^2 \, \nabla^2 \tau_{\mathrm{lat}}(x,t),
\end{equation}
up to errors of order $O(a^4)$. Substituting into~\eqref{eq:discrete-update}
and dividing by $\Delta t$ gives
\begin{equation}
\partial_t \tau_{\mathrm{lat}}
  = \gamma \, \frac{a^2}{\Delta t} \, \nabla^2 \tau_{\mathrm{lat}}
    + \frac{\alpha}{\Delta t} (\rho - \rho_0).
\end{equation}

We define the effective diffusion constant and source coupling
\begin{equation}
D := \gamma \frac{a^2}{\Delta t}, \qquad
\tilde{\alpha} := \frac{\alpha}{\Delta t},
\end{equation}
so that in continuum form
\begin{equation}
\partial_t \tau_{\mathrm{lat}} = D \, \nabla^2 \tau_{\mathrm{lat}} + \tilde{\alpha} (\rho - \rho_0).
\end{equation}

In a stationary configuration ($\partial_t \tau_{\mathrm{lat}} = 0$), this
reduces to
\begin{equation}
D \, \nabla^2 \tau_{\mathrm{lat}} + \tilde{\alpha} (\rho - \rho_0) = 0,
\end{equation}
or
\begin{equation}
\nabla^2 \tau_{\mathrm{lat}} = - \, \frac{\tilde{\alpha}}{D} \, (\rho - \rho_0).
\label{eq:tau-poisson}
\end{equation}

\paragraph{SymPy verification.}
The finite-difference-to-PDE scaling $D = \gamma a^2 / \Delta t$ was checked
symbolically using SymPy in
\\texttt{verification\_notebooks/kappa\_toy\_sympy.ipynb}, confirming that the
6-point stencil converges to $\nabla^2$ with the expected prefactor.

\subsection{Relating Potential $\varphi$ to Latency}
At the interface layer we define the gravitational potential via
\begin{equation}
\varphi = \kappa \, \tau_{\mathrm{lat}}.
\end{equation}
Applying $\nabla^2$ to both sides and using~\eqref{eq:tau-poisson} yields
\begin{equation}
\nabla^2 \varphi = \kappa \, \nabla^2 \tau_{\mathrm{lat}}
  = -\kappa \, \frac{\tilde{\alpha}}{D} \, (\rho - \rho_0).
\end{equation}

To match the Newtonian Poisson equation
\begin{equation}
\nabla^2 \varphi = 4\pi G \, \rho,
\end{equation}
we absorb the $\rho_0$ term into a redefinition of the zero of potential and
require that the coefficient of $\rho$ is $4\pi G$:
\begin{equation}
-\kappa \, \frac{\tilde{\alpha}}{D} = 4\pi G.
\end{equation}
Thus
\begin{equation}
\kappa = -4\pi G \, \frac{D}{\tilde{\alpha}}.
\label{eq:kappa-D-alpha}
\end{equation}

At this stage $\kappa$ is expressed entirely in terms of the diffusion
constant $D$ and the source coupling $\tilde{\alpha}$, which in turn depend
on the microscopic update parameters $\gamma, \alpha, a, \Delta t$.

\subsection{Matching to Queueing-Theoretic Latency}
We now demand consistency between the continuum PDE description of
$\tau_{\mathrm{lat}}$ and the queueing-theoretic expression derived earlier.
In the linear-response regime, queueing theory gave
\begin{equation}
\delta \tau_{\mathrm{lat}} \approx \tau_0 \, \frac{\beta a^3}{C} \, \rho,
\qquad \tau_0 = \frac{f_U \Delta t}{C}.
\label{eq:queue-linear}
\end{equation}

On the other hand, equation~\eqref{eq:tau-poisson} implies that perturbations
in $\tau_{\mathrm{lat}}$ satisfy a Poisson equation with source strength
$-\tilde{\alpha}/D$. For small, nearly homogeneous perturbations one obtains a
linear relationship of the form
\begin{equation}
\delta \tau_{\mathrm{lat}} \propto -\frac{\tilde{\alpha}}{D} \, \rho.
\end{equation}

To make these descriptions compatible, we equate the proportionality factors
up to a choice of normalization that can be absorbed into $\alpha$. For
simplicity we take this normalization to be unity, yielding the matching
condition
\begin{equation}
\frac{\tilde{\alpha}}{D} = -\tau_0 \, \frac{\beta a^3}{C}.
\label{eq:match}
\end{equation}
Substituting $\tau_0 = f_U \Delta t / C$ gives
\begin{equation}
\frac{\tilde{\alpha}}{D} = -\frac{f_U \Delta t}{C} \, \frac{\beta a^3}{C}.
\end{equation}
Using $D = \gamma a^2 / \Delta t$ and $\tilde{\alpha} = \alpha / \Delta t$,
we have
\begin{equation}
\frac{\alpha / \Delta t}{\gamma a^2 / \Delta t}
  = -\frac{f_U \Delta t}{C} \, \frac{\beta a^3}{C},
\end{equation}
so
\begin{equation}
\alpha = -\gamma \, \beta \, f_U \, \Delta t \, \frac{a^5}{C^2}.
\end{equation}

\paragraph{Z3 consistency check.}
The algebraic relations among $D, \tilde{\alpha}, \alpha, \gamma, a, \Delta t$
and the matching condition~\eqref{eq:match} were encoded as real-valued
constraints in Z3 (see
\\texttt{verification\_artifacts/kappa\_constraints\_z3.py}). Z3 reports the
system as satisfiable, confirming internal consistency of the definitions.

\subsection{Final Expression for $\kappa$}
Substituting the matched values of $D$ and $\tilde{\alpha}$ into
~\eqref{eq:kappa-D-alpha} yields
\begin{equation}
\kappa = -4\pi G \, \frac{D}{\tilde{\alpha}} = -4\pi G \, \frac{\gamma a^2}{\alpha}.
\end{equation}
Using the expression for $\alpha$ just found, we obtain
\begin{align}
\kappa &= -4\pi G \, \frac{\gamma a^2}{-\gamma \, \beta \, f_U \, \Delta t \, a^5 / C^2} \\
       &= 4\pi G \, \frac{C^2}{\beta f_U \Delta t} \, \frac{1}{a^3}.
\end{align}
Thus in this toy model
\begin{equation}
\boxed{\kappa
  = 4\pi G \, \frac{C^2}{\beta f_U \Delta t} \, \frac{1}{a^3}.}
\end{equation}

\paragraph{Interpretation.}
Within this explicit substrate model, $\kappa$ is \emph{not} a free parameter.
It is fixed by:
\begin{itemize}
  \item the lattice spacing $a$ (granularity of the substrate),
  \item the update interval $\Delta t$,
  \item the per-site capacity $C$ (bits per time step),
  \item the density-to-arrival coupling $\beta$,
  \item the Universal Clock Frequency $f_U$, and
  \item the empirical gravitational constant $G$ (used here to match the
        interface Poisson equation).
\end{itemize}

This demonstrates, in a fully specified and tool-checked toy model, that
$\kappa$ can emerge as a derived quantity from substrate dynamics rather than
being an arbitrary bridge constant. Generalizing this derivation from a cubic
lattice to the Leech lattice $\leech$ and from a simple diffusion model to the
full UOS dynamics remains an open research direction.

\paragraph{Sage finite-difference validation.}
The convergence of the 6-point discrete Laplacian on a finite cubic lattice to
the continuum Laplacian (in terms of low-lying eigenvalues) was numerically
validated using Sage in
\\texttt{verification\_notebooks/kappa\_lattice\_sage.sage}, providing
additional confidence in the continuum approximation used above.

